\section{Purpose}

\textbf{[STIPULATION]} This document develops the complete scale hierarchy of FFGFT:
from the single geometric constant $\xi$ to the characteristic energy scales, to
the lepton masses, to the fine-structure constant $\alpha$, and to $E_0$ as the
logarithmic mean. The core: all these quantities follow from $\xi$ -- but the
connection is not always achievable by direct cancellation. Ratios (mass ratios,
$E_0$) are correction-free; absolute values need $K_\text{frak}$.

\section{The Starting Point: A Single Constant}

\textbf{[B]} FFGFT begins with a single dimensionless constant:
\[
  \xi = \frac{4}{3}\times10^{-4} = 1{.}333\ldots\times10^{-4}.
\]
From it all fundamental scales follow in natural units ($\hbar=c=\ell_P=1$).

\textbf{[B]} \textbf{Characteristic T0 scales:}
\[
  L_0 = \xi\cdot\ell_P \approx 1{.}333\times10^{-4}\,\ell_P,
  \qquad
  E_0^\text{scale} = L_0^{-1} = \xi^{-1}\cdot E_P.
\]

\section{Lepton Masses from Geometry}

\textbf{[K]} The mass formulas (A100) in compact form:
\[
  m_e = c_e\cdot\xi^{5/2},
  \quad
  m_\mu = c_\mu\cdot\xi^2,
  \quad
  m_\tau = c_\tau\cdot\xi^{3/2},
\]
where the coefficients have geometric origin:
\[
  c_e = \frac{3\sqrt3}{2\pi\alpha^{1/2}},
  \qquad
  c_\mu = \frac{9}{4\pi\alpha},
  \qquad
  c_\tau = \frac{27\sqrt3}{8\pi\alpha^{3/2}}.
\]

\textbf{[K]} \textbf{Important: apparently simple fractions.} In the short form
the coefficients appear as rational numbers ($r_e=4/3$, $r_\mu=16/5$,
$r_\tau=25/9$ from A100). These fractions are \emph{not} arbitrary and must not
be directly cancelled: they encode spatial geometry ($\pi$, $\sqrt3$) and
electromagnetic coupling ($\alpha$). Both representations (Yukawa short form and
extended form) are the same theory in two notations (A130).

\section{$E_0$ as Logarithmic Mean}

\textbf{[B]} The characteristic energy $E_0$ is the geometric mean of electron
and muon mass:
\[
  \boxed{E_0 = \sqrt{m_e\cdot m_\mu}}.
\]
Logarithmically: $\log E_0 = \tfrac12[\log m_e+\log m_\mu]$ -- $E_0$ lies exactly
in the middle between $m_e$ and $m_\mu$ on a logarithmic scale. Numerically:
$E_0 = \sqrt{0{.}511\times105{.}658}\,\text{MeV}\approx7{.}35\,\text{MeV}$
(geometric mean). The value from the $\xi$-path (A130) gives $E_0\approx7{.}398\,\text{MeV}$;
this difference is where $K_\text{frak}$ enters.

\textbf{[B]} From the mass formulas:
\[
  E_0 = \sqrt{c_ec_\mu}\cdot\xi^{9/4}.
\]
$E_0$ is \emph{correction-free} in the sense: because $K_\text{frak}$ appears as
a common factor in both masses, it does not cancel in the geometric mean --
$E_0^\text{exp}=K_\text{frak}\cdot E_0^\text{bare}$, so $E_0$ itself is an
observable that already contains the fractal correction.

\section{The Fine-Structure Constant from $\xi$}

\textbf{[K]} The fundamental relation:
\[
  \alpha = \xi\cdot E_0^2 = c_ec_\mu\cdot\xi^{11/2}.
\]
Substituting the geometric coefficients:
\[
  \alpha^{5/2} = \frac{27\sqrt3}{8\pi^2}\cdot\xi^{11/2},
  \qquad\text{hence}\qquad
  \alpha = \Bigl(\frac{27\sqrt3}{8\pi^2}\Bigr)^{2/5}\cdot\xi^{11/5}.
\]
With fractal correction (A040):
\[
  \alpha = \Bigl(\frac{27\sqrt3}{8\pi^2}\Bigr)^{2/5}\cdot\xi^{11/5}\cdot K_\text{frak},
  \qquad K_\text{frak}=1-100\xi\approx0{.}9867.
\]
Numerically: $\alpha^{-1}\approx137{.}036$ -- agreement to $0{.}0005\,\%$ (A130).

\section{Why $\alpha$ Does Not Follow Directly from $\xi^{11/2}$ Without Unit Convention}

\textbf{[S]} $\alpha\propto\xi^{11/2}$ means: if $\xi$ changed, $\alpha$ would
change with exponent $11/2=5{.}5$. That is a strong coupling. The coefficients
$c_e,c_\mu$ are of order $10^{19}$ in natural units -- that is not a free
parameter but the conversion factor between the Planck scale and MeV. In natural
units (Planck energies) $c_e\sim9{.}67$ and $c_\mu\sim98{.}1$ -- order 1 to 100,
geometrically understandable.

\section{Ratios Are Correction-Free}

\textbf{[B]} The mass ratio:
\[
  \frac{m_\mu}{m_e}
  = \frac{K_\text{frak}\cdot m_\mu^\text{bare}}{K_\text{frak}\cdot m_e^\text{bare}}
  = \frac{m_\mu^\text{bare}}{m_e^\text{bare}}.
\]
$K_\text{frak}$ cancels. The same holds for the Koide ratios (A110) and the
$g$-2 ratios (A138). This is no coincidence: multiplicative corrections always
cancel in ratios.

\textbf{[K]} Absolute values, by contrast, need $K_\text{frak}$: without
correction $\alpha^{-1}\approx138{.}9$ instead of $137{.}04$ -- $1{.}4\,\%$
deviation. With $K_\text{frak}$: exact. Hence $K_\text{frak}$ is not a
post-correction but part of the structure (A040).

\section{The Complete Derivation Chain}

\textbf{[B]} All fundamental quantities from $\xi$ in one chain:
\[
  \xi
  \;\xrightarrow{\text{geometry}}\; r_0,E_0^\text{scale}
  \;\xrightarrow{\text{quantum numbers}}\; m_e,m_\mu,m_\tau
  \;\xrightarrow{\text{geom.\ mean}}\; E_0
  \;\xrightarrow{\xi\cdot E_0^2}\; \alpha
  \;\xrightarrow{\text{K-corr.}}\; \alpha_\text{exp}.
\]
Additionally: $\xi\to L_0\to G$ (gravitational constant, A142) and $\xi\to L_0$ (A015).

\section{Verification Script}

\begin{itemize}[leftmargin=2em,itemsep=2pt,topsep=2pt]
  \item Numerically: $E_0=\sqrt{m_em_\mu}$, $\alpha=\xi E_0^2$, ratios
    correction-free vs.\ absolute values with $K_\text{frak}$:
    \texttt{python/A\_Serie\_Skripte/a261\_skalenhierarchie.py}
\end{itemize}

\section{Open Questions}

\textbf{The coefficients $c_e,c_\mu,c_\tau$ are not fully independently derived.}
They contain $\alpha$ explicitly -- and $\alpha$ in turn follows from the masses.
The circularity is named as a methodological problem in A130; the T0 resolution
(common geometric origin) is a claim, not a proof.

\textbf{The quark coefficients are missing.} The $c_i$ structure for quarks is
not worked out at the same level as for leptons (A150).

\section{Supersedes}

This document subsumes: Dok.~070 (mathematical structure of T0 theory: circular
ratios, scale hierarchy from $\xi$, lepton masses, $\alpha$ derivation, $E_0$ as
logarithmic mean, geometric constant $C$, why ratios are correction-free).
Dok.~093 ($\beta$-parameter derivation, self-consistency condition).

\section{Use of AI Tools}

This document was created with the assistance of AI tools. Responsibility
for content and errors rests with the author. Full wording of the rules in A010.
